\section{Discussion}
\label{sec:discussion}

\subsection{Parser Compliance and Structural Resemblance}

Most outputs parse (255/258), consistent with prior studies reporting high syntax validity~\cite{karpurapu2024,fonseca2025,ferreira2025}; the ablation shows this does not depend on retaining the exemplar. Structural resemblance varies by project: Sylius is above the reference point, while Fineract and Diaspora are below. Long, data-driven Fineract files are often truncated by the encoder, but the 84 reverse-engineered stories based on prefixes do not explain the project difference. These factors were not isolated experimentally. Cross-study similarity scores are contextual, not rankings, because the datasets, metrics and tests differ~\cite{rathnayake2026,fernandes2025}.

\subsection{What Similarity Can and Cannot Establish}

Skeleton cosine measures keyword structure, not behavior. In the exploratory audit, 25 of 60 high-scoring files contained a reviewed defect; cosine did not distinguish defective files, and the scenario-level perturbation test left scores at or above 0.8995. Masked values cannot affect the score by construction. Thus neither 0.85 nor another point in the tested 0.80--0.90 range certifies acceptance-criterion correctness. The weak-to-moderate within-project overlap correlations are consistent with some lexical leakage from reverse-engineered stories, but do not establish causation.

\subsection{Practical Use}

The evidence supports assisted drafting: a reviewer should check expected outcomes, conditions, data values and criterion coverage. The exploratory, project-stratified audit found reviewed defects in 16/20 Fineract files, compared with 6/20 Sylius and 3/20 Diaspora files; this does not estimate defect prevalence in the full corpus. Parser success and high similarity do not support autonomous acceptance or unreviewed CI use.
